\section{Алгоритм численного решения задачи обтекания цилиндра}

В данном разделе приведено пошаговое описание численного алгоритма, реализованного для моделирования дозвукового обтекания цилиндра. Алгоритм базируется на итерационном методе последовательной верхней релаксации по линиям (SLOR).

\subsection{Входные данные и параметры расчета}

Перед началом расчета задаются следующие физические и численные параметры:
\begin{itemize}
    \item $M_\infty$ --- число Маха набегающего потока;
    \item $\gamma = 1.4$ --- показатель адиабаты для воздуха;
    \item $\omega$ --- параметр верхней релаксации (подбирается экспериментально, обычно $1.5 < \omega < 1.8$);
    \item $\epsilon$ --- требуемая точность расчета (критерий остановки по максимальной невязке).
\end{itemize}

\subsection{Генерация расчетной сетки}

Для эффективного решения задачи используется конформное отображение, переводящее внешность цилиндра в полуполосу в плоскости криволинейных координат $(\xi, \eta)$.

\begin{enumerate}
    \item Задается число узлов по радиальной координате $N_\xi$ и по углу $N_\eta$.
    \item Вычисляются шаги сетки: $h_\xi = \xi_{max} / N_\xi$, $h_\eta = 2\pi / N_\eta$.
    \item Связь с физическими координатами $(x, y)$ осуществляется по формулам:
    \[ x_{i,j} = e^{\xi_i} \cos \eta_j, \quad y_{i,j} = e^{\xi_i} \sin \eta_j \]
\end{enumerate}
Такое преобразование позволяет получить сетку, узлы которой сгущаются по экспоненциальному закону при приближении к поверхности цилиндра ($\xi=0$), что критически важно для корректного разрешения больших градиентов параметров.

\subsection{Начальное приближение и граничные условия}

В качестве начального распределения потенциала $\varphi^0$ используется аналитическое решение для обтекания цилиндра несжимаемой жидкостью:
\[ \varphi_{i,j}^0 = (e^{\xi_i} + e^{-\xi_i})\cos\eta_j \]

Граничные условия ставятся следующим образом:
\begin{itemize}
    \item \textbf{На внешней границе} ($\xi = \xi_{max}$): фиксируется потенциал невозмущенного потока (условие Дирихле):
    \[ \varphi_{N_\xi, j} = e^{\xi_{max}} \cos \eta_j \]
    \item \textbf{На поверхности цилиндра} ($\xi = 0$): ставится условие непротекания (условие Неймана):
    \[ \frac{\partial \varphi}{\partial \xi} = 0 \]
    \item \textbf{На линиях $\eta=0$ и $\eta=2\pi$}: используются условия симметрии потока.
\end{itemize}

\subsection{Итерационный процесс решения}

Расчет продолжается до тех пор, пока не будет выполнено условие сходимости:
\[ \max_{i,j} \left| \varphi_{i,j}^{k+1} - \varphi_{i,j}^k \right| < \epsilon \]

Каждая итерация включает в себя два основных этапа:

\subsubsection{Обновление поля плотности}
На основе значений потенциала с предыдущего шага вычисляются компоненты скорости и плотность в каждом узле:
\begin{enumerate}
    \item Производные потенциала в узлах:
    \[ \varphi_\xi \approx \frac{\varphi_{i+1,j} - \varphi_{i-1,j}}{2h_\xi}, \quad \varphi_\eta \approx \frac{\varphi_{i,j+1} - \varphi_{i,j-1}}{2h_\eta} \]
    \item Квадрат модуля скорости:
    \[ q_{i,j}^2 = e^{-2\xi_i} \left( \varphi_\xi^2 + \varphi_\eta^2 \right) \]
    \item Локальная плотность $\rho$ по закону Бернулли для сжимаемого газа:
    \[ \rho_{i,j} = \left[1 + \frac{\gamma-1}{2} M_\infty^2 (1 - q_{i,j}^2)\right]^{\frac{1}{\gamma-1}} \]
\end{enumerate}

\subsubsection{Решение уравнения для потенциала (SLOR)}
Уравнение неразрывности решается итерационным методом по радиальным линиям (кольцам). Для каждой линии $j = \text{const}$ решается трехдиагональная система уравнений 

\[A_i \varphi_{i-1,j}^* + B_i \varphi_{i,j}^* + C_i \varphi_{i+1,j}^* = D_i\].

Коэффициенты матрицы вычисляются с использованием плотностей на гранях ячеек:
\[ A_i = -\frac{\rho_{i-1/2, j}}{h_\xi^2}, \quad C_i = -\frac{\rho_{i+1/2, j}}{h_\xi^2} \]
\[ B_i = \frac{\rho_{i-1/2, j} + \rho_{i+1/2, j}}{h_\xi^2} + \frac{\rho_{i, j-1/2} + \rho_{i, j+1/2}}{h_\eta^2} \]
\[ D_i = \frac{\rho_{i, j-1/2} \varphi_{i, j-1}^{k+1} + \rho_{i, j+1/2} \varphi_{i, j+1}^k}{h_\eta^2} \]

После решения системы методом прогонки применяется процедура верхней релаксации:
\[ \varphi_{i,j}^{k+1} = \varphi_{i,j}^k + \omega \left( \varphi_{i,j}^* - \varphi_{i,j}^k \right) \]
